\documentclass[11pt,a4paper]{article}
% Shared preamble of the FEBAUTO Racing manuals.
\usepackage[margin=2.2cm,headheight=16pt]{geometry}
\usepackage[T1]{fontenc}
\usepackage{lmodern}
\usepackage{xcolor}
\usepackage{graphicx}
\usepackage{booktabs}
\usepackage{tabularx}
\usepackage{listings}
\usepackage{fancyhdr}
\usepackage[hidelinks,colorlinks=true,linkcolor=febblue,urlcolor=febblue]{hyperref}

\definecolor{febblue}{HTML}{1C4687}
\definecolor{febgold}{HTML}{FDB515}
\definecolor{codebg}{HTML}{F3F4F6}

\lstset{
  basicstyle=\ttfamily\small,
  backgroundcolor=\color{codebg},
  frame=single, framerule=0pt, framesep=4pt,
  xleftmargin=4pt, xrightmargin=4pt,
  breaklines=true, columns=fullflexible, keepspaces=true,
  showstringspaces=false,
}
\newcommand{\file}[1]{\texttt{#1}}
\newcommand{\cmd}[1]{\texttt{#1}}
\newcommand{\rulebox}[1]{\par\noindent\colorbox{codebg}{\parbox{\dimexpr\linewidth-2\fboxsep}{#1}}\par\medskip}

\pagestyle{fancy}
\fancyhf{}
\fancyhead[L]{\color{febblue}\textbf{FEBAUTO Racing}}
\fancyhead[R]{\color{febblue}\nouppercase{\leftmark}}
\fancyfoot[C]{\thepage}
\renewcommand{\headrulewidth}{0.8pt}
\renewcommand{\headrule}{\color{febgold}\hrule width\headwidth height\headrulewidth}
\setcounter{tocdepth}{1}

% \febtitle{subtitle}{part lines}
\newcommand{\febtitle}[2]{%
\begin{titlepage}
  \centering
  \vspace*{2cm}
  \includegraphics[width=5cm]{../../assets/feb_logo_512.png}\\[1.5cm]
  {\color{febblue}\Huge\bfseries FEBAUTO Racing}\\[6pt]
  {\Large #1}\\[2cm]
  \begin{tabular}{rl}#2\end{tabular}\\[3cm]
  Formula Electric at Berkeley\\
  September 2026, platform release v0.1.0\\[6pt]
  \url{https://github.com/Pranman1/feb-racing}\\
  \url{https://pranman1.github.io/feb-racing/}
\end{titlepage}
\tableofcontents
\newpage}

\usepackage{array}

\begin{document}
\emergencystretch=2em
\febtitle{Lab: ROS~2 onboarding in the FEB devkit}{%
  \textbf{Session 1} & The devkit, the workspace, packages, nodes, topics, services\\[3pt]
  \textbf{Session 2} & A publisher/subscriber pair, a turtlesim controller, a patrol service\\[3pt]
  \textbf{Source} & \texttt{docs/lab-ros-onboarding.md}; starter code in \texttt{labs/onboarding/}\\}

\section*{About this lab}

One lab, two sessions, adapted from EECS~C106A Labs~1 and~2 (Spring~2026) to run in our devkit
instead of the course's instructional accounts and distrobox container. Every command below was
run as written, in order, on 2026-10-08 on the \texttt{gui} devkit image. Do them in order. The
commands are the same on Mac, Windows and Linux: everything ROS happens inside the container,
and the launcher fills in the ports and paths for your machine.

By the end you can:
\begin{itemize}
  \item start the devkit, open a shell in it and see its desktop in your browser;
  \item create a workspace package with its dependencies, build it with \texttt{colcon}, run nodes with \texttt{ros2 run};
  \item inspect nodes, topics and services with the \texttt{ros2} tools and \texttt{rqt\_graph};
  \item write a publisher/subscriber pair with a custom message type;
  \item write a service server and client, and a node that drives an existing one (turtlesim).
\end{itemize}

Starter code is in \texttt{labs/onboarding/} of the repo: \texttt{chatter} (a talker and a
listener), \texttt{turtle\_patrol\_interface} (the \texttt{Patrol} service definition) and
\texttt{turtle\_patrol} (its server and client). Nothing else is needed. You type the rest.

\section{Setup, once}

You need Docker Desktop (Mac, Windows) or Docker Engine (Linux), git, and about 6~GB of disk.

\begin{lstlisting}
git clone https://github.com/Pranman1/feb-racing.git && cd feb-racing
./feb-sim setup --team "Your Name"
\end{lstlisting}

Already have the clone? \texttt{./feb-sim update} instead: it pulls the repo (the \texttt{labs/}
folder arrived on 2026-10-08), newer devkit images and the newer simulator app in one go.

\texttt{setup} checks Docker, pulls the devkit image and downloads the simulator app. The
simulator is not used in this lab. On a Mac, keep the clone under your home folder
(\texttt{/Users/...}): Docker can only mount folders from there.

\section{What is ROS~2, in one page}

ROS~2 is a set of libraries and tools for robot software. A robot's software is split into small
programs called \textbf{nodes} (a lidar driver, a planner, a motor controller), each doing one
job. Nodes find each other over the network with no central server and talk in three ways:

\begin{itemize}
  \item \textbf{Topics}: a node \emph{publishes} messages on a named topic, any number of nodes
  \emph{subscribe}. A stream, asynchronous, like a broadcast. The lidar publishes \texttt{/scan};
  the planner subscribes.
  \item \textbf{Services}: one node \emph{calls} another and waits for the answer. One request, one
  response, like ordering at a counter. ``Spawn a turtle at (2,~2).''
  \item \textbf{Parameters}: a node's settings, readable and writable at run time.
\end{itemize}

A \textbf{message} is a typed data structure (a \texttt{String}, a \texttt{Twist} with linear and
angular velocity). Nodes can be written in Python (\texttt{rclpy}) or C++ (\texttt{rclcpp}) and
mix freely. Our racing stack is exactly this: a bridge node publishes the simulator's lidar and
encoders as topics, the racer node subscribes, and it publishes steering and throttle topics
back.

\section{The devkit: our container}

The course runs ROS~2 in a distrobox container on the lab machines. We run it in a Docker
container on your laptop. Same idea, the ``sealed lunchbox'': Ubuntu, ROS~2 Humble, all our
dependencies and tools are inside the image; your code stays on your laptop and is mounted in.
Close the container and your laptop is untouched; start it again and it is identical.

Start it with the turtle simulator and a desktop in your browser:

\begin{lstlisting}
./feb-sim ros --gui --apps turtlesim --no-rviz
\end{lstlisting}

What happens: the launcher starts the \texttt{gui} image (first time: a 2.7~GB download), mounts
your \texttt{stack/} folder into it, builds whatever packages are there, starts a screen inside
the container and serves it as a web page, opens \texttt{turtlesim} on that screen, and then
drops you into a shell \textbf{inside the container}. Your browser opens
\url{http://localhost:6080/vnc.html?autoconnect=1&resize=scale} with a turtle in a blue window.
If the tab did not open, paste that URL.

Your prompt now reads \texttt{root@<id>:/home/autodrive\_devkit\#}. You are inside. The
commands you will use from your laptop's own terminal, from the repo folder:

\noindent\begin{tabularx}{\linewidth}{>{\ttfamily\small}p{8.3cm}X}
\toprule
\normalfont Command & What it does \\
\midrule
./feb-sim ros -{}-gui -{}-apps turtlesim -{}-no-rviz & start the devkit with a desktop and turtlesim, shell inside \\
./feb-sim ros -{}-gui -{}-apps "turtlesim rqt" & the same with rviz2 and rqt open too \\
./feb-sim ros & no desktop, no apps: a plain ROS~2 shell (the base image, smaller) \\
./feb-sim shell & another terminal inside the running container; open as many as you like \\
./feb-sim logs -f & what the container printed: the package build, every node it started \\
./feb-sim stop & stop and remove the container; your \texttt{stack/} folder stays \\
./feb-sim ros -{}-src \textasciitilde{}/some/folder & mount a different folder instead of \texttt{stack/} \\
./feb-sim run -{}-track loop & later: the simulator plus the devkit, for the racing labs \\
\bottomrule
\end{tabularx}

\smallskip
You never type a \texttt{docker} command. Each of those is one \texttt{docker run} or
\texttt{docker exec} that the launcher fills in with the right ports and the right folder for
your machine; \texttt{docs/devkit.md} shows the underlying commands if you are curious.

Two habits:
\begin{itemize}
  \item A terminal inside the container is one you opened with \texttt{./feb-sim ros} or
  \texttt{./feb-sim shell}. \texttt{ros2} and \texttt{colcon} only exist there. If a command says
  \texttt{ros2: command not found}, you are in a laptop terminal.
  \item \texttt{exit} in a container shell leaves the container running; \texttt{./feb-sim stop}
  ends it. Do that when you finish, or the next \texttt{./feb-sim ros} replaces it anyway.
\end{itemize}

\section{The workspace}

Inside the container, run \texttt{pwd} and \texttt{ls}:

\begin{lstlisting}
/home/autodrive_devkit
build  desktop.sh  install  log  rviz  src
\end{lstlisting}

This folder is the ROS workspace: the place \texttt{colcon} builds from. It is the same layout
the course lab describes:

\begin{itemize}
  \item \texttt{src/} holds the source of every package. \texttt{src/stack} \textbf{is your laptop's
  \texttt{feb-racing/stack/} folder}, mounted, not copied: a file you save on the laptop is there
  at once, and a package you create inside the container appears on the laptop at once.
  \texttt{src/autodrive\_devkit} and \texttt{src/feb\_tools} are the simulator bridge and house
  tools, part of the image.
  \item \texttt{build/} is where packages are compiled: the prep station.
  \item \texttt{install/} is the result: executables, message code, and the \texttt{setup.bash} that
  tells ROS where it all is. Sourcing it is how ROS finds \emph{your} packages.
  \item \texttt{log/} is what every build printed, in case something fails.
\end{itemize}

\texttt{build/}, \texttt{install/} and \texttt{log/} live in the container and vanish on
\texttt{stop}; that is fine, the container rebuilds \texttt{src/} every time it starts. Only
\texttt{src/stack} is yours and persistent.

\rulebox{\textbf{Linux only:} the container runs as root, so a package created \emph{inside} it is
owned by root on your laptop and your editor cannot save it. Once, after creating packages, run
on the laptop: \texttt{sudo chown -R \$USER stack}. Mac and Windows Docker Desktop map the
ownership for you.}

\section{Anatomy of a package}

A \textbf{package} is the unit of ROS code: a folder with a \texttt{package.xml} describing it
and its dependencies, plus source, launch files, message definitions. Look at the one we will
play with:

\begin{lstlisting}
cd /opt/ros/humble/share/turtlesim
ls
cat package.xml
\end{lstlisting}

\texttt{package.xml} names the package, its maintainers and licence, and what it depends on:
\texttt{<build\_depend>} at build time, \texttt{<exec\_depend>} at run time,
\texttt{<depend>} both. \texttt{rclcpp} there is the C++ client library; Python packages depend on \texttt{rclpy}
instead, and that is what we use.

Tools that answer questions about packages (\texttt{-{}-help} on any of them lists the options):

\begin{lstlisting}
ros2 pkg prefix turtlesim                          # where a package is installed
ros2 pkg executables turtlesim                     # the programs it provides
ros2 interface show example_interfaces/msg/String  # the fields of a message type
ros2 pkg prefix autodrive_roboracer                # our simulator bridge, built into the image
\end{lstlisting}

\section{Creating and building packages}

Every package you make goes in \texttt{src/stack} (anywhere else is inside the container's own
disk and lost on stop). Create two: a bare one and one with dependencies declared.

\begin{lstlisting}
cd /home/autodrive_devkit/src/stack
ros2 pkg create foo --build-type ament_python
ros2 pkg create bar --build-type ament_python --dependencies rclpy std_msgs geometry_msgs turtlesim
ls foo
grep depend bar/package.xml
\end{lstlisting}

\texttt{ament} is ROS~2's build system; \texttt{ament\_python} is a pure Python package,
\texttt{ament\_cmake} a C++ one (or one that only holds message definitions, see
section~\ref{sec:custom}). In \texttt{foo} you find \texttt{package.xml} (metadata,
dependencies), \texttt{setup.py} (how to install the Python code, and where executables are
registered), \texttt{setup.cfg}, \texttt{resource/foo} (an empty marker file), \texttt{foo/} (the
Python module your nodes go in, with an empty \texttt{\_\_init\_\_.py}) and \texttt{test/} (style
checks, ignore them). Open them on your laptop: they are at
\texttt{feb-racing/stack/foo}.

Build from the workspace root, not from inside the package:

\begin{lstlisting}
cd /home/autodrive_devkit
echo $AMENT_PREFIX_PATH
colcon build --symlink-install --packages-select foo bar
source install/setup.bash
echo $AMENT_PREFIX_PATH
ros2 pkg prefix bar
\end{lstlisting}

\texttt{-{}-packages-select} builds just the ones you name (a plain \texttt{colcon build} also
rebuilds the bridge, which works but takes longer). \texttt{-{}-symlink-install} links the
installed Python files to your source, so an edit takes effect on the next \texttt{ros2 run}
without rebuilding; you rebuild only when you add a package, an executable or a message.
\texttt{source install/setup.bash} adds your packages to the paths ROS searches, which is the
difference between the two \texttt{echo} lines. Every new shell needs it (\texttt{./feb-sim
shell} does it for you).

\rulebox{\textbf{Checkpoint 1.} Be able to explain the workspace layout and what \texttt{build},
\texttt{install}, \texttt{log} and \texttt{src/stack} are; show the build; explain
\texttt{package.xml}; find a package's path.}

\section{Nodes, topics and services}

turtlesim is already running on the desktop (section~3). Ask ROS what is there:

\begin{lstlisting}
ros2 node list
ros2 node info /turtlesim
\end{lstlisting}

\texttt{/turtlesim} is the node. It subscribes to \texttt{/turtle1/cmd\_vel} (velocity
commands), publishes \texttt{/turtle1/pose}, and serves \texttt{/spawn}, \texttt{/clear},
\texttt{/kill} and more. (\texttt{/foxglove\_bridge} is the devkit's own node that lets the
Foxglove app on your laptop see topics; ignore it.)

Drive the turtle. In a \textbf{second} terminal on your laptop:

\begin{lstlisting}
./feb-sim shell
ros2 run turtlesim turtle_teleop_key
\end{lstlisting}

Click into that terminal and use the arrow keys: the turtle in the browser moves. This is
\texttt{ros2 run <package> <executable>}: the teleop node reads your keys and publishes
\texttt{Twist} messages on \texttt{/turtle1/cmd\_vel}; the turtlesim node subscribes and moves
the turtle.

See the graph. Leave the teleop running, and in a third terminal (\texttt{./feb-sim shell}):

\begin{lstlisting}
rqt_graph &
\end{lstlisting}

It opens on the browser desktop next to the turtle (drag windows by their title bar; the
desktop is 1600 by 900). Two ovals, one arrow, one topic. Under \emph{Hide}, untick \emph{Debug}
and refresh: the terminals listening in appear too.

Topics by hand:

\begin{lstlisting}
ros2 topic list
ros2 topic list -v                       # with message types and counts
ros2 topic type /turtle1/cmd_vel
ros2 topic echo /turtle1/pose            # Ctrl+C to stop; drive the turtle and watch
ros2 topic echo /turtle1/cmd_vel
\end{lstlisting}

Services are the other pattern: one request, one response.

\begin{lstlisting}
ros2 service list
ros2 service type /clear
ros2 service call /clear std_srvs/srv/Empty "{}"
\end{lstlisting}

The turtle's track is wiped. \texttt{/spawn} takes arguments; find out which:

\begin{lstlisting}
ros2 service type /spawn
cat $(ros2 pkg prefix turtlesim)/share/turtlesim/srv/Spawn.srv
\end{lstlisting}

Above the \texttt{-{}--} is the request (x, y, theta, name), below it the response (the name
given).

\begin{lstlisting}
ros2 service call /spawn turtlesim/srv/Spawn "{x: 2.0, y: 2.0, theta: 1.2, name: 'new_turtle'}"
\end{lstlisting}

A second turtle appears, and \texttt{ros2 topic list} now has \texttt{/new\_turtle/...} topics.

\rulebox{\textbf{Checkpoint 2.} Explain node, topic, message, service; drive the turtle; show a
topic's type and echo it; spawn a turtle.}

\section{A publisher and a subscriber}

The starter package \texttt{chatter} is in the repo under \texttt{labs/}, which is not mounted
into the container, so copy it on your \textbf{laptop}, in the repo folder:

\begin{lstlisting}
cp -r labs/onboarding/chatter stack/
\end{lstlisting}

Then inside the container:

\begin{lstlisting}
cd /home/autodrive_devkit
colcon build --symlink-install --packages-select chatter
source install/setup.bash
ros2 pkg executables chatter
\end{lstlisting}

Two executables, \texttt{talker} and \texttt{listener}. Run the first in this terminal and the
second in another (\texttt{./feb-sim shell}):

\begin{lstlisting}
ros2 run chatter talker
ros2 run chatter listener
\end{lstlisting}

The talker prints \texttt{Publishing: "Hello World: n"} twice a second; the listener prints
\texttt{I heard:} with the same text. Start a second listener: both hear every message. Start a
second talker: the listener hears both, interleaved. Ctrl+C stops each.

Read \path{stack/chatter/chatter/publisher_member_function.py} and
\path{subscriber_member_function.py} on your laptop. A node is a class that inherits
\texttt{rclpy.node.Node}; \path{create_publisher(type, topic, depth)} and
\path{create_subscription(type, topic, callback, depth)} make the ends;
\texttt{create\_timer(seconds, callback)} runs something regularly; \texttt{rclpy.spin(node)}
keeps the node alive, running its timers and callbacks until Ctrl+C.

Now \texttt{stack/chatter/setup.py}, at the bottom:

\begin{lstlisting}
entry_points={
    'console_scripts': [
        'talker = chatter.publisher_member_function:main',
        'listener = chatter.subscriber_member_function:main',
    ],
},
\end{lstlisting}

This is what makes \texttt{ros2 run chatter talker} exist: executable name, then
\texttt{module:function}. The file can be called anything; the executable name is what you
declare here. Without an entry point there is nothing for \texttt{ros2 run} to run.

\section{Write your own pair, with a custom message}
\label{sec:custom}

Build a talker that publishes what you type, with a timestamp, and a listener that prints it
with the time it arrived. Read this whole section first.

\textbf{The message type.} \texttt{std\_msgs/String} has one field; you need two (text and a
timestamp). Custom messages are inconvenient in an \texttt{ament\_python} package, so the
convention is a separate \texttt{ament\_cmake} \emph{interface package} that holds only message
and service definitions:

\begin{lstlisting}
cd /home/autodrive_devkit/src/stack
ros2 pkg create my_chatter_msgs --build-type ament_cmake
mkdir my_chatter_msgs/msg
\end{lstlisting}

Create \path{my_chatter_msgs/msg/TimestampString.msg} with two lines, one field each: a
\texttt{string} and an \texttt{int64}. Then tell the build about it. In
\texttt{my\_chatter\_msgs/package.xml}, after the existing \texttt{<buildtool\_depend>} line, add:

\begin{lstlisting}
<buildtool_depend>rosidl_default_generators</buildtool_depend>
<exec_depend>rosidl_default_runtime</exec_depend>
<member_of_group>rosidl_interface_packages</member_of_group>
\end{lstlisting}

In \texttt{my\_chatter\_msgs/CMakeLists.txt}, before \texttt{ament\_package()}, add:

\begin{lstlisting}
find_package(rosidl_default_generators REQUIRED)
rosidl_generate_interfaces(${PROJECT_NAME}
  "msg/TimestampString.msg"
)
find_package(rosidl_default_runtime REQUIRED)
ament_export_dependencies(rosidl_default_runtime)
\end{lstlisting}

Build and check ROS can see it:

\begin{lstlisting}
cd /home/autodrive_devkit
colcon build --symlink-install --packages-select my_chatter_msgs
source install/setup.bash
ros2 interface show my_chatter_msgs/msg/TimestampString
\end{lstlisting}

\textbf{The nodes.} Create \texttt{my\_chatter} (\texttt{ament\_python}, dependencies
\texttt{rclpy my\_chatter\_msgs}) in \texttt{src/stack}, and write two files in its module
folder, using \texttt{chatter} as your model:

\begin{itemize}
  \item \texttt{my\_talker.py}: in a loop, \texttt{input()} a line from the terminal, make a
  \texttt{TimestampString} with the text and \path{self.get_clock().now().nanoseconds},
  publish it on \texttt{/user\_messages}. Import the type with
  \path{from my_chatter_msgs.msg import TimestampString}.
  \item \texttt{my\_listener.py}: subscribe to \texttt{/user\_messages} and print
  \path{Message: <text>, Sent at: <stamp>, Received at: <stamp>}, where the second stamp
  is the listener's own clock when the callback ran. It is not in the message; think about why
  it can differ from the first.
\end{itemize}

Add entry points \texttt{my\_talker} and \texttt{my\_listener} to \texttt{setup.py}, build with
\texttt{-{}-packages-select my\_chatter}, source, and run each in its own \texttt{./feb-sim shell}:

\begin{lstlisting}
ros2 run my_chatter my_talker
ros2 run my_chatter my_listener
\end{lstlisting}

Type lines into the talker; they appear in the listener with both stamps.

\rulebox{\textbf{Checkpoint 3.} Explain everything in \texttt{stack/}; show the message definition
and that both packages build; demonstrate the pair.}

\section{A controller for turtlesim (Lab 2, part 1)}
\label{sec:controller}

Replace \texttt{turtle\_teleop\_key} with your own node. Create \texttt{lab2\_turtlesim}
(\texttt{ament\_python}, dependencies \texttt{rclpy std\_msgs geometry\_msgs turtlesim}). Its
node \texttt{turtle\_controller}:

\begin{enumerate}
  \item Takes the turtle's name on the command line:
\begin{lstlisting}
ros2 run lab2_turtlesim turtle_controller <name_of_turtle>
\end{lstlisting}
  Use \texttt{sys.argv}, after \path{rclpy.utilities.remove_ros_args(sys.argv)} strips ROS's
  own flags.
  \item Publishes \texttt{geometry\_msgs/msg/Twist} on that turtle's \texttt{cmd\_vel} topic when
  the user presses keys. Reading single keystrokes from a terminal is fiddly; \texttt{input()}
  plus Enter is fine (\texttt{w}/\texttt{a}/\texttt{s}/\texttt{d}, say). Bonus if you capture
  keys directly.
  \item Works for several turtles at once: spawn more with \texttt{/spawn} (section~7) and run one
  controller per turtle in its own \texttt{./feb-sim shell}, each with a different name.
\end{enumerate}

One thing you will meet: a node that publishes within a few hundred milliseconds of starting
can lose those first messages, because turtlesim has not yet discovered the new publisher.
Typing is slower than that, so it does not matter here, but remember it when a node's first
message ``disappears''.

\rulebox{\textbf{Checkpoint 4.} Explain the package; show it builds; drive two turtles from two
controllers.}

\section{Services: a patrol server (Lab 2, part 2)}

A \textbf{service} is a request and a response between two nodes: a client sends a request, the
server's callback runs and fills in a response, the client gets it. For something you ask for
once, on demand; topics are for streams.

Three parts define one: a service type (\texttt{.srv} file), a server node, a client node. The
starter has all three. On the laptop:

\begin{lstlisting}
cp -r labs/onboarding/turtle_patrol_interface labs/onboarding/turtle_patrol stack/
\end{lstlisting}

Read \path{stack/turtle_patrol_interface/srv/Patrol.srv}:

\begin{lstlisting}
float32 vel
float32 omega
---
geometry_msgs/Twist cmd
\end{lstlisting}

Request above the line, response below. Note in its \texttt{CMakeLists.txt} the
\texttt{DEPENDENCIES geometry\_msgs} line: a service that uses another package's message must say
so, and \texttt{package.xml} must depend on that package.

\path{stack/turtle_patrol/turtle_patrol/patrol_server.py} offers \texttt{/turtle1/patrol}:
a request stores a \texttt{Twist} with the requested speed and turn rate, and a timer publishes
it on \texttt{/turtle1/cmd\_vel} ten times a second, so the turtle circles.
\texttt{patrol\_client.py} calls it once: \texttt{create\_client}, \texttt{wait\_for\_service},
\texttt{call\_async}, then \texttt{spin\_until\_future\_complete} to wait for the answer. Build
and run (server and client in separate shells):

\begin{lstlisting}
cd /home/autodrive_devkit
colcon build --symlink-install --packages-select turtle_patrol_interface turtle_patrol
source install/setup.bash
ros2 run turtle_patrol patrol_server
ros2 run turtle_patrol patrol_client
\end{lstlisting}

turtle1 circles. The service can also be called from the command line, with other values:

\begin{lstlisting}
ros2 service call /turtle1/patrol turtle_patrol_interface/srv/Patrol "{vel: 2.0, omega: 1.0}"
\end{lstlisting}

\textbf{Your extension: one server, many turtles, with teleport.} Change the system so that a
single server controls any turtle by name, teleports it to a requested pose first, then patrols
it:

\begin{lstlisting}
ros2 service call /turtle_patrol turtle_patrol_interface/srv/Patrol \
  "{turtle_name: 'new_turtle', vel: 2.0, omega: 1.0, x: 1.0, y: 3.0, theta: 1.0}"
\end{lstlisting}

Requirements:
\begin{enumerate}
  \item One server process, \texttt{multi\_patrol\_server}, serving \texttt{/turtle\_patrol}. Not
  one server per turtle. The request's \texttt{turtle\_name} says which turtle.
  \item The request also carries \texttt{x}, \texttt{y}, \texttt{theta}. The server teleports that
  turtle there with turtlesim's own teleport service, then starts its patrol with \texttt{vel}
  and \texttt{omega}.
  \item A client, \texttt{patrol\_client}, called as
  \path{ros2 run turtle_patrol patrol_client <name> <x> <y> <theta> <vel> <omega>}, that
  fills the request from those arguments, calls the service and prints the response.
\end{enumerate}

What to change: the \texttt{.srv} (add \texttt{string turtle\_name}, \texttt{float32 x},
\texttt{float32 y}, \texttt{float32 theta} to the request; add to the response, for instance
\texttt{bool success} and \texttt{string message}, but keep \texttt{geometry\_msgs/Twist cmd} or
the starter server and client stop working), then rebuild the interface package before the
Python one.

Hints:
\begin{itemize}
  \item Find the teleport service with \path{ros2 service list}, then its type with
  \path{ros2 service type <name>} and its fields with \path{ros2 interface show <type>}.
  Your server is a \emph{client} of it; model that part on \texttt{patrol\_client.py}.
  \item Keep per-turtle state in a dictionary keyed by name: the publisher for
  \texttt{/<name>/cmd\_vel} and the current \texttt{Twist}. The timer loops over the dictionary.
  \item Do not wait for the teleport inside your service callback with
  \texttt{spin\_until\_future\_complete}: the node is already spinning to run your callback, and
  waiting there deadlocks. \texttt{call\_async} and carry on; the teleport completes a moment
  later.
\end{itemize}

\rulebox{\textbf{Checkpoint 5.} Explain how a service call differs from a subscriber callback;
show two turtles patrolling from different starting poses; explain why one server is enough.}

\section{Keep your code in git}

\texttt{stack/} is ignored by the \texttt{feb-racing} repo on purpose: it is yours. Make it a
repo of its own and push it to a \textbf{private} GitHub repository.

Once, set up a key (skip if GitHub already knows this laptop):

\begin{lstlisting}
git config --global user.name "Your Name"
git config --global user.email you@berkeley.edu
ssh-keygen -t ed25519 -C "you@berkeley.edu"       # Enter for the default file; a passphrase is wise
cat ~/.ssh/id_ed25519.pub                           # paste into GitHub -> Settings -> SSH and GPG keys
ssh -T git@github.com                               # "Hi <you>! You've successfully authenticated"
\end{lstlisting}

Create an empty private repository on github.com, then:

\begin{lstlisting}
cd stack
git init
git remote add origin git@github.com:<you>/<repo>.git
git add .
git commit -m "ROS onboarding lab"
git push -u origin main
\end{lstlisting}

From then on: \texttt{git add}, \texttt{git commit -m "..."}, \texttt{git push}. Partners may
share the repository; it stays private.

\section{When something is wrong}

\noindent\begin{tabularx}{\linewidth}{lX}
\toprule
Symptom & Cause, fix \\
\midrule
\texttt{ros2: command not found} & you are in a laptop terminal; \texttt{./feb-sim shell} \\
\texttt{Package 'x' not found} after building & this shell has not sourced the new build: \texttt{source install/setup.bash} \\
\texttt{ros2 pkg create} made files you cannot edit (Linux) & root-owned: \texttt{sudo chown -R \$USER stack} on the laptop \\
the browser desktop says not connected & reload the tab; still nothing, \texttt{./feb-sim logs -n 50} \\
\texttt{mounts denied} (Mac) & the clone is outside \texttt{/Users}: move it \\
a package I created is gone after \texttt{stop} & it was created outside \texttt{src/stack}; recreate it there \\
my node's first message never arrives & discovery: see section~\ref{sec:controller} \\
the build fails & \texttt{./feb-sim logs -n 100}, or run \texttt{colcon build} again and read the error \\
\bottomrule
\end{tabularx}

\section*{Checkoff summary}

\begin{enumerate}
  \item Workspace, \texttt{colcon build}, \texttt{package.xml}, \texttt{ros2 pkg prefix}.
  \item Node, topic, message, service; teleop, echo, spawn.
  \item \texttt{my\_chatter\_msgs} and \texttt{my\_chatter} built and demonstrated.
  \item \texttt{lab2\_turtlesim} driving two turtles.
  \item \texttt{multi\_patrol\_server} patrolling two turtles from requested poses; code pushed.
\end{enumerate}

\end{document}
